\subsection{同伦扩张性质}

定义 8 —— 设 X 是拓扑空间，A 是 X 的一个子集。若对每个拓扑空间 Y、每个连续映射 \(f:\mathrm{X}\to\mathrm{Y}\) 以及每个同伦 \(\sigma:\mathrm{A}\times\mathbf{I}\to\mathrm{Y}\)（其终点为映射 \(f|_{\mathrm{A}}\)），都存在同伦 \(\tau:\mathrm{X}\times\mathbf{I}\to\mathrm{Y}\)（它扩张 \(\sigma\)，且终点为映射 f），则称偶 (X,A) 具有同伦扩张性质。

注 1 —— 设 X 是拓扑空间，A 是 X 的子集，且偶 (X,A) 具有同伦扩张性质。设 Y 是拓扑空间，\(f:\mathrm{X}\to\mathrm{Y}\) 是连续映射，\(\sigma:\mathrm{A}\times\mathrm{I}\to\mathrm{Y}\) 是起点为映射 \(f|_{\mathrm{A}}\) 的同伦。映射 \(\overline{\sigma}:\mathrm{A}\times\mathrm{I}\to\mathrm{Y}\)（由 \((a,t)\mapsto\sigma(a,1-t)\) 定义）是终点为 \(f|_{\mathrm{A}}\) 的同伦；设 \(\tau:\mathrm{X}\times\mathrm{I}\to\mathrm{Y}\) 是扩张 \(\overline{\sigma}\) 且终点为 f 的同伦。于是映射 \(\overline{\tau}:\mathrm{X}\times\mathrm{I}\to\mathrm{Y}\)（由 \((x,t)\mapsto\tau(x,1-t)\) 给出）是扩张 \(\sigma\) 且起点为映射 f 的同伦。

设 X 是拓扑空间，A 是 X 的子空间，\(i: A \to X\) 是典范单射。用 \(\alpha_{i}: A \times I \to \text{Cyl}(i)\) 与 \(\beta_{i}: X \to \text{Cyl}(i)\) 表示典范映射。设 \(j: \text{Cyl}(i) \to X \times I\) 是唯一的连续映射，使得有 \(j(\alpha_{i}(a, s)) = (i(a), s)\) 与 \(j(\beta_{i}(x)) = (x, 1)\)（对 \(a \in A\)、\(s \in I\) 及 \(x \in X\)）。它是单射；其像是子空间 \((\mathbf{A} \times \mathbf{I}) \cup (\mathbf{X} \times \{1\})\)（\(X \times I\) 的）。若 A 在 X 中闭，则映射 j 是闭的。它并不总是严格的（III, p. 325, 习题 17）。

命题 7 —— 沿用上述记号，下列断言等价：

\begin{enumerate}
\def\labelenumi{(\roman{enumi})}
\item
  偶 (X, A) 具有同伦扩张性质；
\item
  对每个拓扑空间 Y 及每个连续映射 \(g\colon \operatorname{Cyl}(i) \to \operatorname{Y}\)，存在连续映射 \(g'\colon X \times I \to Y\)，使得 \(g = g' \circ j\)；
\item
  单射 j 有一个连续的收缩映射；
\item
  映射 j 是严格的，且存在 X × I 到 j 的像上的一个收缩。
\end{enumerate}

设偶 (X,A) 具有同伦扩张性质。设 \(g:\operatorname{Cyl}(i)\to Y\) 是连续映射；令 \(\sigma=g\circ\alpha_i\)，\(f=g\circ\beta_i\)。映射 \(f:X\to Y\) 是连续的，而 \(\sigma:A\times I\to Y\) 是终点为映射 \(f|A\) 的同伦。因此存在同伦 \(g':X\times I\to Y\)（终点为 f，且扩张 \(\sigma\)）。我们有 \(g'\circ j(\alpha_i(a,s))=g'(a,s)=\sigma(a,s)=g(\alpha_i(a,s))\) 与 \(g'\circ j(\beta_i(x))=g'(x,1)=f(x)=g(\beta_i(x))\)（对 \(a\in A\)、\(s\in I\) 及 \(x\in X\)）。于是 \(g'\circ j=g\)，从而得到 (ii)。

反之，设断言 (ii) 成立，我们来证明偶 (X,A) 具有同伦扩张性质。设 Y 是拓扑空间，\(f:X\to Y\) 是连续映射，\(\sigma:A\times I\to Y\) 是终点等于 \(f|A\) 的同伦。存在唯一的连续映射 \(g:\mathrm{Cyl}(i)\to\mathrm{Y}\)，使得 \(g\circ\alpha_i=\sigma\) 且 \(g\circ\beta_i=f\)（III, p. 238, 命题 4）。于是，映射 \(g':\mathrm{X}\times\mathbf{I}\to\mathrm{Y}\) 若满足 \(g=g'\circ j\)，便是终点为 f 且扩张 \(\sigma\) 的同伦。

  断言 (ii) 与 (iii) 的等价性是 III, p. 235 引理 1 的特殊情形。

设单射 \(j\colon \mathrm{Cyl}(i)\to \mathrm{X}\times \mathbf{I}\) 有一个连续收缩映射 \(r\)。映射 \(j\) 定义从 \(\mathrm{Cyl}(i)\) 到子空间 \(\mathrm{T} = (\mathrm{A}\times \mathbf{I})\cup (\mathrm{X}\times \{1\})\)（\(\mathrm{X}\times \mathbf{I}\) 的）上的同胚（出处同上）。因此它是严格的。

令 \(\rho = \mathrm{pr}_1\circ j\circ r\)，\(\theta = \mathrm{pr}_2\circ j\circ r\)。若 \(x\in \mathrm{A}\) 或 \(s = 1\)，我们有 \(\rho (x,s) = x\) 与 \(\theta (x,s) = s\)。对 \((x,s)\in \mathrm{X}\times \mathbf{I}\) 及 \(t\in \mathbf{I}\)，令

\[
\sigma ((x, s), t) = (\rho (x, (1 - t) + s t), (1 - t) s + t \theta (x, s)).
\]

映射 \(\sigma\colon(\mathrm{X}\times\mathbf{I})\times\mathbf{I}\to\mathrm{X}\times\mathbf{I}\) 是连续的。对 \((x,s)\in\mathrm{X}\times\mathbf{I}\)，我们有

\[
\sigma ((x, s), 0) = (\rho (x, 1), s) = (x, s)
\]

以及

\[
\sigma ((x, s), 1) = (\rho (x, s), \theta (x, s)) = j \circ r (x, s);
\]

对 \(x\in \mathbf{X}\) 及 \(t\in \mathbf{I}\)，我们有

\[
\sigma ((x, 1), t) = (\rho (x, 1), (1 - t) + t \theta (x, 1)) = (x, 1),
\]

而对 \((x,s)\in\mathbf{A}\times\mathbf{I}\) 及 \(t\in\mathbf{I}\)，我们有

\[
\sigma ((x, s), t) = (x, (1 - t) s + t s) = (x, s).
\]

因此，\(\sigma\) 是 \(X\times I\) 到 T 上的一个收缩。这表明断言 (iii) 蕴涵断言 (iv)。

蕴涵关系 (iv)⇒(iii) 是显然的。

注 2 —— 设 X 是分离拓扑空间，A 是 X 的子集，且偶 (X,A) 具有同伦扩张性质。用 \(i\colon A\to X\) 表示典范单射。设 \(j\colon\operatorname{Cyl}(i)\to X\times\mathbf{I}\) 是典范单射，\(r\) 是映射 j 的一个连续收缩映射。子空间 \(j(\operatorname{Cyl}(i))\) 等于偶 \((x,t)\in X\times\mathbf{I}\)（满足 \(j(r(x,t))=(x,t)\)）所成的集合。由于空间 \(X\times\mathbf{I}\) 是分离的，子集 \(j(\operatorname{Cyl}(i))\) 在 \(X\times\mathbf{I}\) 中是闭的。于是集合 A——它等于 \(x\in X\)（使 \((x,0)\in j(\operatorname{Cyl}(i))\)）所成的集合——是 X 的闭子集。

引理 3 —— 设 X 与 Y 是拓扑空间，\(p \colon X \to Y\) 是逆紧且开的连续映射，\(f \colon X \to \mathbf{R}\) 是连续映射。映射 \(g \colon Y \to \overline{\mathbf{R}}\)（由 \(y \mapsto \sup_{x\in p^{-1}(y)} f(x)\) 给出）是连续的。

设 \(b \in \mathbf{Y}\)。由于 \(p\) 逆紧，其纤维 \(\overline{p}^1(b)\) 是拟紧空间；因此我们有 \(g(b) \in \mathbf{R} \cup \{-\infty\}\)。

设 \(m \in R\) 满足 \(g(b) < m\)。对所有 \(a \in \overline{p}^{1}(b)\)，我们有 \(f(a) \leqslant g(b) < m\)；设 \(V_{a}\) 是 a 在 X 中的邻域，使得 \(f(x) < m\)（对所有 \(x \in V_{a}\)）。诸集合 \(V_{a}\) 的并 V 是 \(\overline{p}^{1}(b)\) 在 X 中的邻域。由引理 5（I, p. 75），存在 b 在 Y 中的邻域 W，使得 \(\overline{p}^{1}(W) \subset V\)。对所有 \(y \in W\)，我们有 \(g(y) \leqslant m\)。这证明 g 在 b 处是上半连续的。

现在证明 \(g\) 在 \(b\) 处是下半连续的。我们可以假设 \(g(b)\in\mathbf{R}\)。设 \(m\in\mathbf{R}\) 满足 \(m<g(b)\)。设 \(a\in\overline{p}^1(b)\) 满足 \(m<f(b)\)；于是设 V 是 a 在 X 中的邻域，使得 \(f(x)>m\)（对所有 \(x\in V\)）。由此得出 \(g(y)>m\)（对所有 \(y\in p(V)\)）。由于 p 是开的，\(p(V)\) 是 b 的邻域，因此 g 是下半连续的。

引理至此证毕。

定理 1 —— 设 X 是拓扑空间，A 是 X 的闭子空间；用 i : A → X 表示典范单射。下列断言等价：

\begin{enumerate}
\def\labelenumi{(\roman{enumi})}
\item
  偶 (X, A) 具有同伦扩张性质；
\item
  存在连续映射 \(\varphi\colon\mathrm{X}\to\mathbf{I}\)，使得 \(\mathrm{A}=\overline{\varphi}^{1}(0)\)，并存在在 A 上固定、终点为 X 的恒同映射的同伦 \(\sigma\colon\mathrm{X}\times\mathbf{I}\to\mathrm{X}\)，使得有 \(\sigma(x,0)\in\mathrm{A}\)（对每个点 \(x\in\mathrm{X}\)，满足 \(\varphi(x)\neq1\)）；
\item
  存在连续映射 \(\varphi\colon\mathbf{X}\to\mathbf{R}_{+}\)，使得 \(\mathrm{A}=\overline{\varphi}^{1}(0)\)，并存在在 A 上固定的同伦 \(\sigma\colon\overline{\varphi}^{1}(\mathbf{I})\times\mathbf{I}\to\mathrm{X}\)，使得 \(\sigma(x,1)=x\) 且 \(\sigma(x,0)\in\mathrm{A}\)（对所有 \(x\in\overline{\varphi}^{1}(\mathbf{I})\)）；
\item
  存在连续映射 \(\varphi\colon X\to\mathbf{R}_{+}\)，使得 \(A\subset\overline{\varphi}^{1}(0)\)，并存在连续映射
\end{enumerate}

\[
\sigma \colon \{(x, t) \in \mathrm{X} \times \mathbf {I} \mid t + \varphi (x) \geqslant 1 \} \to \mathrm{X}
\]

使得有 \(\sigma(x,1) = x\)（对所有 \(x \in X\)）且 \(\sigma(x,1 - \varphi(x)) \in A\)（对所有 \(x \in X\)，满足 \(\varphi(x) \leqslant 1\)）。

我们用 i 表示 A 到 X 中的典范单射，用 \(\alpha_{i}\colon A\times I\to\operatorname{Cyl}(i)\) 与 \(\beta_{i}\colon X\to\operatorname{Cyl}(i)\) 表示典范映射，并用 j 表示从 \(\operatorname{Cyl}(i)\) 到 \(X\times I\) 的映射，它满足：\(j(\alpha_{i}(x,s))=(x,s)\)（对 \((x,s)\in A\times I\)），\(j(\beta_{i}(x))=(x,1)\)（对 \(x\in X\)）。

设偶 (X,A) 具有同伦扩张性质，并设 \(r:X\times I\to\operatorname{Cyl}(i)\) 是映射 j 的一个连续收缩映射（III, p. 241, 命题 7）。用 \(\sigma\) 表示连续映射 \(\operatorname{pr}_1\circ j\circ r\)（从 \(X\times I\) 到 X）。对每个 \((x,t)\in X\times I\)，我们有 \(|\operatorname{pr}_2(j(r(x,t)))-t|\leqslant1\)。由于 I 是紧的，第一投影 \(\operatorname{pr}_1:X\times I\to X\) 是逆紧映射；它也是开的。由引理 3，映射 \(\varphi:X\to I\)（由

\[
x \mapsto \sup _ {t \in \mathbf {I}} \left(\left| \operatorname{pr} _ {2} (j (r (x, t))) - t \right|\right)
\]

给出）是连续的。

对 \(x \in \mathrm{X}\)，我们有 \(\sigma(x,1) = x\)。设 \(x \in \mathrm{X}\) 满足 \(\sigma(x,0) \notin \mathrm{A}\)；则 \(\mathrm{pr}_2(j(r(x,0))) = 1\) 且 \(\varphi(x) \geqslant 1\)。对 \(x \in \mathrm{A}\) 及 \(t \in \mathbf{I}\)，我们有 \(j(r(x,t)) = (x,t)\)；因此 \(\sigma(x,t) = x\) 且 \(\varphi(x) = 0\)。反之，设 \(x \in \mathrm{X}\) 满足 \(\varphi(x) = 0\)。于是有 \(\mathrm{pr}_2(j(r(x,t)) \leqslant t\)（对所有 \(t \in \mathbf{I}\)）；若 \(t < 1\)，这蕴涵 \(j(r(x,t)) \in \mathrm{A} \times \mathbf{I}\)；由于 A 在 X 中闭，我们因此有 \(j(r(x,1)) \in \mathrm{A} \times \mathbf{I}\)，从而 \(x \in \mathrm{A}\)。

这表明 (i) 蕴涵 (ii)。

设 \(\varphi\) 与 \(\sigma\) 是满足断言 (ii) 中性质的映射。令 \(\varphi_1 = 2\varphi\)，并设 \(\sigma_1\) 是 \(\sigma\) 在 \(\overline{\varphi}_1^1 (\mathbf{I})\times \mathbf{I}\) 上的限制。我们有 \(\overline{\varphi}_1^1 (0) = \mathrm{A}\)；对 \(a\in \mathrm{A}\)，有 \(\sigma_1(a,t) = a\)（对所有 \(a\in \mathrm{A}\)）。设 \(x\in \mathrm{X}\) 满足 \(\varphi_{1}(x)\leqslant 1\)；我们有 \(\sigma_1(x,1) = x\)；此外，\(\varphi (x) = \varphi_{1}(x) / 2 < 1\)，因此 \(\sigma_{1}(x,0)\in \mathrm{A}\)。于是 (ii) 蕴涵 (iii)。

我们来证明 (iii) 蕴涵 (iv)。设 \(\varphi\colon X\to\mathbf{R}\) 与 \(\sigma\colon\overline{\varphi}^{1}(\mathbf{I})\times\mathbf{I}\to X\) 如命题所述；令 \(B=\overline{\varphi}^{1}(\mathbf{I})\)，\(C=\overline{\varphi}^{1}([1,+\infty[)\)。

设 \(u_1\) 是从 B × I 到 X 的映射：\(u_1(x,t) = \sigma(x,1 - (1 - t)/2\varphi(x))\)（若 \(t + 2\varphi(x) \geqslant 1\) 且 \(\varphi(x) > 0\)），否则 \(u_1(x,t) = \sigma(x,0)\)。由下面的引理 4，它是连续的。

设 \(u_2\) 是从 B × I 到 X 的映射：\(u_2(x,t)=\sigma(x,\sup(0,1-2(1-t)(1-\varphi(x))))\)。它是连续的。对 \(x\in\mathrm{B}\)（满足 \(\varphi(x)=1/2\)）及 \(t\in\mathbf{I}\)，我们有 \(u_1(x,t)=\sigma(x,t)=u_2(x,t)\)。用 \(u:\mathrm{B}\times\mathbf{I}\to\mathrm{X}\) 表示如下的映射：\(u(x,t)=u_1(x,t)\)（若 \(\varphi(x)\leqslant1/2\)），\(u(x,t)=u_2(x,t)\)（若 \(1/2<x\leqslant1\)）。它是连续的，因为它在 B×I 的闭子空间 \(\overline{\varphi}^{1}([0,1/2])\times\mathbf{I}\) 与 \(\overline{\varphi}^{1}([1/2,1])\times\mathbf{I}\) 上的限制都是连续的（TG, I, p. 19, 命题 4）。

对 \(x\in X\)（满足 \(\varphi(x)=1\)）及 \(t\in\mathbf{I}\)，我们有 \(u(x,t)=u_2(x,t)=\sigma(x,1)=x\)。因此存在唯一的映射 \(\tau:X\times I\to X\)，它与 u 在 \(B\times I\) 上一致，与映射 \(\mathrm{pr}_1\) 在 \(C\times I\) 上一致；它是连续的（出处同上）。

设 \(x \in X\)；可以验证 \(\tau(x,1) = 1\)。此外，若 \(x \in A\)，则 \(\varphi(x) = 0\)，因此 \(\tau(x,t) = u_{1}(x,t) = \sigma(x,0) = x\)。最后，若 \(2\varphi(x) \leqslant 1\)，则 \(\varphi(x) \leqslant 1/2\)，因此 \(\tau(x,1 - 2\varphi(x)) = \sigma(x,0) \in A\)。

这就证明了断言 (iv) 成立。

最后，设断言 (iv) 满足，我们来证明偶 (X, A) 具有同伦扩张性质。

用 \(C_{1}\)（相应地，\(C_{2}\)）表示偶 \((x,t)\in\mathbf{X}\times\mathbf{I}\)（满足 \(t+\varphi(x)\leqslant1\)（相应地，\(t+\varphi(x)\geqslant1\)））所成的集合。它们都是闭集。对 \((x,t)\in\mathrm{C}_{1}\)，我们有 \(\sigma(x,1-\varphi(x))\in\mathrm{A}\)；设 \(\rho_{1}\colon\mathrm{C}_{1}\to\mathrm{Cyl}(i)\) 是由 \((x,t)\mapsto\alpha_{i}(\sigma(x,1-\varphi(x)),t+\varphi(x))\) 给出的映射；它是连续的。再设 \(\rho_{2}\colon\mathrm{C}_{2}\to\mathrm{Cyl}(i)\) 是由 \((x,t)\mapsto\beta_{i}(\sigma(x,t))\) 给出的连续映射。对 \((x,t)\in\mathrm{C}_{1}\cap\mathrm{C}_{2}\)，我们有 \(t+\varphi(x)=1\)，因此

\[
\rho_ {1} (x, t) = \alpha_ {i} (\sigma (x, 1 - \varphi (x)), 1) = \beta_ {i} (\sigma (x, 1 - \varphi (x)) = \rho_ {2} (x, t).
\]

因此存在唯一的映射 \(\rho\colon X\times\mathbf{I}\to\mathrm{Cyl}(i)\)，它与 \(\rho_{1}\) 在 \(C_{1}\) 上一致，与 \(\rho_{2}\) 在 \(C_{2}\) 上一致；它是连续的（TG, I, p. 19, 命题 4）。

对 \(x \in A\) 及 \(t \in I\)，我们有 \(\varphi(x) = 0\)，因此 \(t + \varphi(x) \leqslant 1\)。

\[
\rho (j (\alpha_ {i} (x, t))) = \rho (x, t) = \alpha_ {i} (\sigma (x, 1), t) = \alpha_ {i} (x, t).
\]

对 \(x\in \mathrm{X}\)，我们还有

\[
\rho (j (\beta_ {i} (x))) = \rho (x, 1) = \beta_ {i} (\sigma (x, 1)) = \beta_ {i} (x).
\]

由于映射 \(\alpha_{i}\) 与 \(\beta_{i}\) 的像覆盖 \(\mathrm{Cyl}(i)\)，由此得出映射 \(\rho\) 是映射 j 的连续收缩映射。于是，偶 (X, A) 具有同伦扩张性质（III, p. 241, 命题 7），定理证毕。

引理 4 —— 设 X 与 Y 是拓扑空间，\(\varphi\colon\mathrm{X}\to\mathbf{R}_{+}\) 是连续映射，并令 \(\mathrm{A}=\overline{\varphi}^{1}(0)\)。设 \(\sigma\colon\mathrm{X}\times\mathbf{I}\to\mathrm{Y}\) 是在 A 上固定的同伦。映射 \(\sigma'\colon\mathrm{X}\times\mathbf{I}\to\mathrm{Y}\) 把 \((x,s)\) 映到 \(\sigma(x,s/\varphi(x))\)（若 \(s<\varphi(x)\)）且映到 \(\sigma(x,1)\)（若 \(s\geqslant\varphi(x)\)），它是连续的。

映射 \(\sigma'\) 在闭子集 \(\overline{\varphi}^{1}([1, +\infty[)\times\mathbf{I}\) 的每一点处都连续；因此只需证明它在 \(\overline{\varphi}^{1}(\mathbf{I})\times\mathbf{I}\) 上的限制是连续的。于是我们可以假设 \(\varphi(\mathrm{X})\subset\mathbf{I}\)。设 C 与 C' 分别是由偶 \((x,s)\)（满足 \(s\leqslant\varphi(x)\) 与 \(s\geqslant\varphi(x)\)）组成的 X × I 的子空间。它们是闭的且覆盖 X × I。映射 \(\sigma'\) 在 C' 上连续；我们来证明它在 C 上连续。

设 \(\alpha\colon\mathbf{X}\times\mathbf{I}\to\mathbf{X}\times\mathbf{I}\) 是由 \(\alpha(x,t)=(x,t\varphi(x))\) 给出的连续映射。它的像等于 C，且 \(\sigma'\circ\alpha\) 是连续映射 \(\sigma\)。我们来证明 \(\alpha\) 是逆紧映射。事实上，考虑超滤子 \(\mathfrak{U}\)（在 \(\mathbf{X}\times\mathbf{I}\) 上）以及点 \((x,t)\in\mathbf{X}\times\mathbf{I}\)（附着于超滤子基 \(\alpha(\mathfrak{U})\)）。由于 \(\mathrm{pr}_{1}\circ\alpha=\mathrm{pr}_{1}\)，X 上的超滤子基 \(\mathrm{pr}_{1}(\mathfrak{U})\) 收敛于 \(x\)。由于 I 是紧的，存在点 \(s\in\mathbf{I}\) 使超滤子基 \(\mathrm{pr}_{2}(\mathfrak{U})\) 收敛于 \(s\)。于是 \(\mathfrak{U}\) 收敛于 \((x,s)\)。由于 \(\alpha\) 连续，超滤子基 \(\alpha(\mathfrak{U})\) 收敛于 \((x,s\varphi(x))\)。由于 I 是分离的，我们有 \(s\varphi(x)=t\)，于是 \(\alpha(x,s)=(x,t)\)，从而 \(\alpha\) 是逆紧的（TG, I, p. 75, 定理 1）。再由 I, p. 18, 例 2 及命题 9 推出 \(\sigma'\mid\mathrm{C}\) 是连续的。

推论 1 —— 设 X 是正规拓扑空间，A 是 X 的闭子空间。假设存在 A 在 X 中的邻域 V 及 V 到 A 上的一个收缩，以及连续映射 \(f: \mathrm{X} \to \mathbf{R}\)（满足 \(f^{-1}(0) = \mathrm{A}\)）。那么偶 (X,A) 具有同伦扩张性质。

设 \(\rho\colon V \times I \to V\) 是 \(V\) 到 \(A\) 上的一个收缩。由正规空间的定义（TG, IX, p. 41, 定义 1），存在连续映射 \(g\colon X \to I\)，它在 \(A\) 上取值 0，在 \(X - V\) 的每点处取值 1。设 \(\varphi\colon X \to R\) 是由 \(\varphi(x) = |f(x)| + g(x)\)（\(x \in X\)）给出的映射；它是连续的。我们有 \(\overline{\varphi}^{1}(0) = A\) 与 \(\overline{\varphi}^{1}(I) \subset V\)。设 \(\sigma\) 是从 \(\overline{\varphi}^{1}(I) \times I\) 到 \(X\) 的映射，由 \(\sigma(x,t) = \rho(x,1-t)\) 给出（对 \(x \in \overline{\varphi}^{1}(I)\) 及 \(t \in I\)）。对 \(x \in \overline{\varphi}^{1}(I)\)，我们有 \(\sigma(x,1) = \rho(x,0) = x\) 与 \(\sigma(x,0) = \rho(x,1) \in A\)。

映射 \(\varphi\) 与 \(\sigma\) 满足 III, p. 243 定理 1 断言 (iii) 的条件；因此，偶 (X,A) 具有同伦扩张性质。

例 1 —— 取球 \(\mathbf{B}_{n}\) 作为空间 X，球面 \(\mathbf{S}_{n-1}\) 作为子空间 A。若 \(V = X - \{0\}\)，则存在 V 到 \(\mathbf{S}_{n-1}\) 上的一个强收缩（III, p. 237, 例）。因此，偶 \((\mathbf{B}_{n}, \mathbf{S}_{n-1})\) 具有同伦扩张性质。

推论 2 —— 设 X 与 Y 是拓扑空间，A 是 X 的闭子空间，B 是 Y 的闭子空间。若偶 (X,A) 与 (Y,B) 都具有同伦扩张性质，则偶 (X × Y, (X × B) ∪ (A × Y)) 也具有同伦扩张性质。

设 \(\varphi\colon X\to\mathbf{R}_{+}\) 与 \(\sigma\colon\overline{\varphi}^{1}(\mathbf{I})\times\mathbf{I}\to X\)，以及 \(\varphi^{\prime}\colon Y\to\mathbf{R}_{+}\) 与 \(\sigma^{\prime}\colon\overline{\psi}^{1}(\mathbf{I})\times\mathbf{I}\to Y\)，分别是关于偶 (X,A) 与偶 (Y,B) 满足定理 1 断言 (iv) 的条件的映射。设 \(\psi\colon X\times Y\to\mathbf{R}_{+}\) 是由 \((x,y)\mapsto\inf(\varphi(x),\varphi^{\prime}(x))\) 给出的映射；它是连续的；我们还有 \(\psi(x,y)=0\)（对所有 \((x,y)\in X\times Y\)，满足 \(x\in A\) 或 \(y\in B\)）。对 \((x,y,t)\in X\times Y\times I\)（满足 \(t+\psi(x)\geqslant1\)），我们有 \(t+\varphi(x)\geqslant1\) 与 \(t+\varphi^{\prime}(x)\geqslant1\)。于是可以定义连续映射

\[
\tau \colon \{(x, y, t) \in \mathrm{X} \times \mathrm{Y} \times \mathbf {I} \mid t + \psi (x) \geqslant 1 \} \rightarrow \mathrm{X} \times \mathrm{Y}
\]

如下：令 \(\tau(x,y,t) = (\sigma(x,t),\sigma'(y,t))\)。对所有 \((x,y)\in X\times Y\)，我们有 \(\tau(x,y,1) = (\sigma(x,1),\sigma'(y,1)) = (x,y)\)。此外，若 \(\psi(x,y)\leqslant 1\)，则 \(\varphi(x)\leqslant 1\) 或 \(\varphi'(y)\leqslant 1\)，从而 \(\sigma(x,1)\in A\) 或 \(\sigma'(y,1)\in B\)；这蕴涵 \(\tau(x,y,1 - \psi(x))\) 属于 \((A\times Y)\cup (X\times B)\)。这就验证了定理 1 的断言 (iv)，推论得证。

例 2 —— *下面是偶 (X, A) 具有同伦扩张性质的其他情形。

\begin{enumerate}
\def\labelenumi{(\roman{enumi})}
\item
  空间 X 是 C \(^{1}\) 类仿紧微分流形，A 是 X 的闭子流形。鉴于推论 \(^{1}\)，这由如下事实得出：X 是完全正规的（IX, p. 103, 习题 11），且 A 在 X 中有一个管状邻域；
\item
  空间 X 是胞腔复形，A 是相对于 X 的一个给定胞腔分解的子复形。*
\end{enumerate}
% semantic-fragments: legacy-input-seam
\relax{}
